\documentclass{article}

\usepackage[utf8]{inputenc}     % pdfLaTeX
\usepackage[T1]{fontenc}
\usepackage[magyar]{babel}

\usepackage{amsthm}
\usepackage{amsmath}
\usepackage{amsfonts}
\usepackage{tikz}
\usetikzlibrary{arrows.meta, positioning}
\usepackage{graphicx}

\title{Kombinatorikai feladatok megoldása kódolással}
\author{ATaKoMa Studios et al}
\date{Legutóbb frissítve: 2026. Április 24.}

\begin{document}

\maketitle
\tableofcontents
\newpage

\section{Mozgás}

A körödben: \begin{enumerate}
    \item Végrehajthatsz szektorok közötti lépést
    \item Végrehajthatsz szektoron belüli lépést
    \item Passzolhatsz
\end{enumerate}

\subsection{Szektorok közötti lépés}
Egy tanuló csoporttal léphetsz különböző szektorokba akármilyen távolságot, amit ki tudsz fizetni. Minden érintett szektor 1 energiába kerül (szomszédos szektorba menni 1, egy szektoron át menni és vele szomszédosba menni 2 energia...)

\subsection{Szektoron belüli lépés}
Egy vagy két tanulóddal lépj a szektoron belül. Amennyiben olyan terembe lépsz, amelyben már áll tanulód, plusz egy energiát kell fizetned.

Szektoron belüli lépés után a tanulód "lefekszik", ebben a fordulóban már nem mozoghatsz vele

\subsection{Passz}
A saját körödben dönthetsz úgy, hogy passzolsz. Amennyiben passzolsz, ebben a fordulóban már nem kell lépned, de nem is léphetsz többet.

\section{Küldetések}

Amennyiben olyan terembe lépsz be, amelyben van tanár akitől nincs küldetésed, vehetsz fel tőle küldetést. Egy tanulóval akár két különböző nehézségű küldetést is felvehetsz.

A küldetések lehetnek könnyűek, közepesek és nehezek. 

Amennyiben van nálad könnyű küldetés, nem teljesíthetsz közepeset és nehezet.  

Egy tanárnál egyszerre 3 küldetés látszik egy könnyű egy közepes és egy nehéz. Egy tanárnak 14  küldetése van, 4 játékos esetén 4-3-3, 5 játékos esetén 5-4-3, 6 játékos esetén 6-4-4 küldetés van játékban)

Egyszerre egy játékosnál maximum 10 küldetés lehet. 

\section{Tanulók beszerzése}



\end{document}